\subsection{Observable equivalence and sufficient state variables}

A common carrier is useful only together with an observable equivalence
relation. Let $\mathcal X$ be an admitted state domain, $\Phi_t$ its
full evolution, and $\pi:\mathcal X\to\mathcal Y$ a readout.
Two states are equivalent for this readout when $\pi(x)=\pi(x')$.
An induced evolution $\phi_t$ on the image exists precisely when
\begin{equation}
 \pi(x)=\pi(x')\quad\Longrightarrow\quad
 \pi(\Phi_t x)=\pi(\Phi_t x')
\end{equation}
for all admitted states and times. Necessity follows by applying
$\phi_t$ to the common initial readout. For sufficiency, define
$\phi_t(y)=\pi(\Phi_t x)$ using any representative with $\pi(x)=y$;
the condition makes the definition independent of the representative.
For a full flow preserving the domain, its composition law descends
as well. Existence, regularity, and locality of a field flow remain
separate analytic questions.

This elementary criterion connects two seemingly different issues.
A compressed field register can omit orientation needed for a Lorentz
transformation. A projected state can omit an amplitude needed for
the derivative of its phase. In both cases equivalence classes of
the original readout are split by the proposed operation. The remedy
must supply the missing information, restrict the admitted domain,
or describe a family of possible outputs. It cannot be supplied by
renaming the original register.

\subsection{Three tests of a proposed physical interface}

We organize the investigation around three questions. First, what
is represented? This requires a carrier, scalar field, adjoint and
operator action. The distinction between an involution and an internal
complex structure is already a necessary test at this level. Second,
which operation selects the physical states? Closure, positive
spectral selection and a generation seed are different restrictions.
Their intersection may have smaller dimension than any one separately.
Third, which measured quantity is read out? An expectation, a transition
norm, a matrix eigenvalue and a pole of a dressed propagator need not
give the same mass law.

The answers can be checked before a detailed phenomenological fit.
Rank rules out three positive eigenmasses from a one-dimensional
transition output. Positivity bounds exclude a globally positive
linear-density realization at the intended root radius. A hidden-lift
counterexample excludes a universal phase-only evolution for a given
projection. These are useful negative results because their hypotheses
identify which part of a model would have to change. They do not
exclude every representation on the common carrier.

The presentation follows these interfaces rather than the historical
order of research rounds. After the carrier and its comparison metrics,
an electromagnetic example explains the assertion and execution
language. Particle representations are separated from generation
seeds and mixing. The mass section then follows a transition operator
through its Gram matrix, a gapped effective comparison, and the positive
root constraints. The geometric section examines projected evolution,
action reduction, coframes and finite compatibility. Finally, a zeta
comparison tests whether arithmetic and uncertainty can supply an
additional selection operation. Original supporting mass arguments are
retained in Appendix~A; the full Companion \cite{QSC411} supplies the
broader research setting and alternative routes.
